\documentclass{article}

\usepackage{microtype}
\usepackage{graphicx}
\usepackage{booktabs}
\usepackage{hyperref}
\usepackage{url}
\usepackage[accepted]{icml2022}

\icmltitlerunning{Multi-Agent Workflow Design with AFlow}

\begin{document}

\twocolumn[
\icmltitle{Multi-Agent Workflow Design with AFlow for Mathematical and Code Reasoning}

\begin{icmlauthorlist}
\icmlauthor{Student Name}{pku}
\end{icmlauthorlist}

\icmlaffiliation{pku}{Student ID: TODO; Course: Multi-Agent Reinforcement Learning}
\icmlcorrespondingauthor{Student Name}{TODO@example.com}

\icmlkeywords{multi-agent systems, large language models, workflow optimization, AFlow}

\vskip 0.3in
]

\begin{abstract}
This report studies agentic workflow construction for large language model reasoning tasks using the AFlow codebase. The assignment goal is to understand AFlow's engineering structure, reproduce direct and chain-of-thought baselines, and design multi-agent workflows for MATH-500 and HumanEval. We imported AFlow, configured a Kimi K2.5-compatible OpenAI-style client, implemented reproducible experiment runners, and added manual multi-agent workflows for math and code tasks. Twenty-sample validation runs verify the evaluation, API, workflow, token logging, and ablation pipelines end to end. Full benchmark runs remain the next step before drawing final performance conclusions.
\end{abstract}

\section{Introduction}

Large language model agents can improve reasoning by decomposing a task into multiple model calls, critiques, revisions, code execution steps, or ensemble decisions. AFlow formulates this idea as a search problem over code-represented agentic workflows, where reusable operators become nodes in a workflow graph and an optimizer searches for better graph structures \cite{zhang2025aflow}. Related work on multi-agent debate \cite{du2024debate} and optimizable agent graphs \cite{zhuge2024gptswarm} motivates comparing direct model calls against explicit multi-agent reasoning.

The assignment requires experiments on MATH-500 and HumanEval. MATH-style tasks test symbolic and quantitative reasoning \cite{hendrycks2021math}; HumanEval tests executable Python function synthesis \cite{chen2021humaneval}. The current project focuses first on a reproducible engineering base: safe API configuration, benchmark runners, manual workflows, validation subsets, and result tables.

\section{AFlow Engineering Summary}

The imported AFlow snapshot is stored under \texttt{AFlow/}. The main engineering components are:

\begin{itemize}
    \item \textbf{Nodes and LLM calls}: \texttt{scripts/async\_llm.py} wraps OpenAI-compatible async chat completion APIs and tracks token usage.
    \item \textbf{Operators}: \texttt{scripts/operators.py} and dataset-specific operator templates provide reusable calls such as generation, self-consistency ensemble, code generation, execution-based testing, and revision.
    \item \textbf{Workflows}: \texttt{workspace/<Dataset>/workflows/<name>/graph.py} implements callable Python workflow graphs.
    \item \textbf{Benchmarks}: \texttt{benchmarks/math.py} and \texttt{benchmarks/humaneval.py} evaluate model outputs against ground truth or executable tests.
    \item \textbf{Optimization}: \texttt{run.py} wires datasets, operators, and the optimizer for automatic workflow search.
\end{itemize}

\section{Environment Setup}

The environment follows AFlow's Python 3.9 recommendation:

\begin{verbatim}
conda create -n marl_hw2 python=3.9 -y
conda install -n marl_hw2 pip=24.2 -y
conda run -n marl_hw2 python -m pip install -r AFlow/requirements.txt
\end{verbatim}

Two dependency issues were found during setup. First, the imported downloader required \texttt{requests}, which was not listed in the snapshot requirements. Second, the local SOCKS proxy required \texttt{socksio} for \texttt{httpx}. MATH scoring also needed \texttt{antlr4-python3-runtime} for SymPy's LaTeX parser. These were added to \texttt{AFlow/requirements.txt}.

The local Kimi configuration is created from \texttt{AFlow/config/config2.kimi.example.yaml}; \texttt{AFlow/config/config2.yaml} is ignored by Git and reads the secret from \texttt{KIMI\_API\_KEY}. Initial Kimi API smoke tests established the required non-thinking parameters \texttt{temperature=0.6} and \texttt{top\_p=0.95}. A separate \texttt{kimi-k2.5-thinking} alias is kept for model-side thinking experiments.

\section{Baselines}

The project implements reproducible baseline runners in \texttt{experiments/run\_baselines.py}. Each run records the dataset, model, split, sample indices, concurrency, CSV path, and aggregate score in a generated \texttt{run\_config.json}.

\paragraph{Direct baseline.} The direct math baseline asks Kimi K2.5 to solve the problem and place the final answer in \texttt{\textbackslash boxed\{\}}. The direct HumanEval baseline asks for valid Python code defining the required entry point.

\paragraph{CoT baseline.} The chain-of-thought baseline adds explicit step-by-step reasoning instructions. For HumanEval, the prompt allows analysis but asks the final answer to contain valid Python code, which the evaluator sanitizes before execution.

\section{Manual Workflow Design}

\paragraph{MATH manual\_v1.} The math workflow in \texttt{AFlow/workspace/MATH/workflows/manual\_v1/} creates three independent solution candidates with different reasoning instructions. It then applies self-consistency selection and a final cleanup call that enforces a single \texttt{\textbackslash boxed\{\}} answer. The intended effect is to reduce single-sample reasoning errors while preserving compatibility with AFlow's answer extractor.

\paragraph{HumanEval manual\_v1.} The code workflow in \texttt{AFlow/workspace/HumanEval/workflows/manual\_v1/} creates two independent code candidates. If public tests are available, it runs the dataset-specific \texttt{Test} operator and uses reflection repair for failing candidates. If no candidate passes public tests, the workflow falls back to self-consistency selection over the repaired candidates.

\section{Validation Results}

Table~\ref{tab:validation20} reports 20-sample validation subset tests using sample seed 1. These are not final benchmark results, but they are large enough to reveal initial differences between direct prompting, CoT, manual workflows, and ablations. MATH scores in this table are recomputed from saved predictions with the corrected text-answer normalizer, so CSV filenames may still contain the original pre-fix score.

\begin{table}[t]
\caption{Twenty-sample validation subset tests.}
\label{tab:validation20}
\vskip 0.08in
\begin{center}
\begin{small}
\begin{tabular}{llrrr}
\toprule
Dataset & Method & Score & Calls & Tokens \\
\midrule
HumanEval & direct & 1.00000 & 20 & 6,034 \\
HumanEval & CoT & 1.00000 & 20 & 15,628 \\
HumanEval & manual\_v1 & 1.00000 & 41 & 16,441 \\
HumanEval & no public test & 1.00000 & 62 & 38,973 \\
MATH & direct & 0.85000 & 20 & 16,557 \\
MATH & CoT & 0.95000 & 20 & 17,708 \\
MATH & manual\_v1 & 0.95000 & 100 & 151,117 \\
MATH & single candidate & 0.85000 & 40 & 54,945 \\
\bottomrule
\end{tabular}
\end{small}
\end{center}
\end{table}

On MATH, CoT improves over direct prompting from 0.85 to 0.95. The full \texttt{manual\_v1} workflow also reaches 0.95, while the single-candidate ablation remains at 0.85. This suggests that the multi-candidate selection path fixes at least one direct-generation failure, but the current workflow does not yet outperform CoT and costs substantially more tokens. On HumanEval, all methods score 1.00 on this validation subset, so the subset is saturated and cannot distinguish workflow quality. The no-public-test ablation is especially costly because it generates and selects among candidates without early execution-feedback filtering.

Failure analysis found that two MATH misses were evaluator false negatives caused by text-answer formatting such as \texttt{MAKE} versus \texttt{\textbackslash text\{MAKE\}} and comma spacing inside text answers. The evaluator now normalizes simple text wrappers and comma spacing before symbolic comparison. Remaining MATH failures are genuine reasoning/counting issues, such as a rectangle-counting problem and an overinclusive vector-projection answer.

\section{Planned Full Experiments}

The next experimental stage is to expand from 20-sample validation tests to larger controlled subsets and then full benchmark runs. The planned sequence is:

\begin{enumerate}
    \item Run 30--50 validation examples for direct, CoT, and \texttt{manual\_v1} where validation set size permits.
    \item Inspect failure logs to distinguish model mistakes from formatting or extraction errors.
    \item Tune prompts and workflow branching only on validation data.
    \item Run full MATH-500 and HumanEval test evaluations.
    \item Add larger ablations: remove ensemble, remove public-test repair, reduce candidate count, and compare non-thinking against thinking-enabled Kimi K2.5.
\end{enumerate}

\section{Limitations}

Current results are validation subsets and should not be interpreted as final benchmark gains. The manual workflows increase token usage and latency, especially for MATH where three candidates plus ensemble selection used 151,117 tokens on 20 problems. Kimi pricing is not yet reflected in the local cost field because the imported AFlow pricing table does not include Kimi. Future runs should report raw token counts or add a verified Kimi pricing table.

\section{Conclusion}

The project now has a reproducible AFlow-based implementation path: safe Kimi configuration, dataset download, direct and CoT baselines, manual multi-agent workflows, ablations, token-aware validation tables, and a report draft. The remaining core work is to scale experiments further and improve the MATH workflow beyond the CoT baseline before final test evaluation.

\bibliography{references}
\bibliographystyle{icml2022}

\end{document}
